\documentclass[11pt]{article}

\usepackage[margin=1in]{geometry}
\usepackage{amsmath, amssymb}
\usepackage{graphicx}
\usepackage{listings}
\usepackage{xcolor}
\usepackage{hyperref}
\setlength{\parskip}{0.5em}
\setlength{\parindent}{0pt}

\lstset{
    basicstyle=\ttfamily\small,
    breaklines=true,
    showstringspaces=false,
    frame=single,
    language=Python,
}

\title{Project 1: Least Squares \\ \large Oil Prices and U.S. Unemployment}
\author{} % TODO: your name(s)
\date{\today}

\begin{document}

\maketitle

\begin{abstract}
% TODO: 1-2 paragraphs. Overview of experimental question and main conclusions.
\end{abstract}

\section{Introduction}
% TODO: utility/importance of curve fitting; brief statement of experimental question.

\section{Methods and Computations}

\subsection{Dataset}
% TODO: thorough description of the dataset and its provenance (include URLs).
% Source: FRED (Federal Reserve Economic Data)
%   - POILBREUSDM: https://fred.stlouisfed.org/series/POILBREUSDM
%   - UNRATE:      https://fred.stlouisfed.org/series/UNRATE

\subsection{Comparison of Linear Algorithms}
% TODO: recap the math of normal equations + Cholesky, QR + back-substitution,
% and the built-in solver. Compare results with figures/tables.

\subsection{Condition Number}
% TODO: relate cond(A) (and cond(A^T A)) to the stability of each method.

\subsection{Other Curve Fits}
% TODO: quadratic, cubic, or other nonlinear fits; discuss whether they're
% reasonable given the data, with figures/tables.

\section{Summary of Results}
% TODO: discuss findings in relation to the experimental question.

\section{Conclusions}
% TODO: key takeaways; lessons for interpreting curve fits in the news/media.

\appendix
\section{Code}
% TODO: insert full commented code listing, e.g.:
% \lstinputlisting[language=Python]{oil_unemployment_demo.py}

\end{document}
